% -------------------------------------------------------------------------------------------------
%      MDSG Latex Framework
%      ============================================================================================
%      File:                  abstract.tex
%      Author(s):             Michael Duerr
%      Version:               1
%      Creation Date:         30. Mai 2010
%      Creation Date:         30. Mai 2010
%
%      Notes:                 - Place your abstract here
% -------------------------------------------------------------------------------------------------
%
\vspace*{2cm}

\begin{center}
    \textbf{Abstract}
\end{center}

\vspace*{1cm}

\noindent Unsupervised Environment Design (UED) has emerged as a promising paradigm in the single-agent reinforcement learning (RL) setting, where automatically generated curricula of progressively more complex environments lead to robust and generalizable policies.

The goal of this thesis is a systematic transfer of the UED approach to cooperative multi-agent reinforcement learning (MARL) domains. As such, the thesis shall explore the design, implementation and evaluation of environments with generalizing elements - i.e., underspecified parameters - that may be realized as various parameterized aspects of cooperation and coordination. The underspecified parameters of the environment may further include structural constraints such as layout, resource distribution, and partial observability. As an application domain, a cooperative benchmark with parameterizable layouts, for example from the Overcooked domain or similar Grid-based navigation task domains, may be considered. The dynamic aspects of the UED method, the task-design, the environment scenarios and curriculum discovery suitability will be analyzed in detail.

This thesis shall further investigate the extent to which current autocurricula from the literature — such as Prioritized Level Replay (PLR), ACCEL, SFL or MAESTRO — can be applied to parameterized multi-agent environments that require cooperation, in order to induce emergent cooperative behaviors through increasing environmental complexity. For this purpose, conventional reinforcement learning algorithms or heuristics from the literature, as well as other self-generated baselines may be used and compared. The findings aim to contribute to a deeper understanding of how open-ended environment design facilitates the emergence of cooperation in multi-agent systems.